% Mock draft of the final paper, laid out like the course's IEEE template
% (project-template-letter.docx). Build the PDF from the docs/ folder with:
%     pdflatex paper_draft.tex
% Red text marks TODOs and [X] placeholders that must come from real results.

\documentclass[10pt,twocolumn,letterpaper]{article}
\usepackage[margin=0.75in,columnsep=0.25in]{geometry}
\usepackage{mathptmx}          % Times font, like the IEEE template
\usepackage[T1]{fontenc}
\usepackage{amssymb}           % \square for checklist boxes
\usepackage{booktabs}
\usepackage{enumitem}
\usepackage{xcolor}
\usepackage{caption}
\usepackage{titlesec}
\usepackage[hidelinks]{hyperref}

% IEEE-style headings: I. INTRODUCTION, A. Subsection
\renewcommand\thesection{\Roman{section}}
\renewcommand\thesubsection{\Alph{subsection}}
\titleformat{\section}{\centering\normalsize\scshape}{\thesection.}{0.5em}{}
\titleformat{\subsection}{\normalsize\itshape}{\thesubsection.}{0.5em}{}
\titlespacing*{\section}{0pt}{1.2ex}{0.8ex}
\titlespacing*{\subsection}{0pt}{1ex}{0.5ex}
\captionsetup[table]{labelformat=simple,labelsep=newline,font={footnotesize,sc},
  justification=centering}
\renewcommand\thetable{\Roman{table}}
\captionsetup[figure]{font=footnotesize,name=Fig.,labelsep=period}

% Placeholders
\newcommand{\todo}[1]{\textcolor{red}{\textbf{TODO:} #1}}
\newcommand{\X}{\textcolor{red}{[X]}}
\newcommand{\figbox}[1]{\fbox{\parbox[c][0.95in][c]{0.92\columnwidth}{\centering
  \textcolor{red}{\small Figure goes here:\\ \texttt{#1}}}}}

\title{\Large Depth versus Strength: Evaluating Minimax with\\ Alpha-Beta Pruning in Connect Four\\[4pt]
  \normalsize\textcolor{red}{MOCK DRAFT --- not for submission}}
\author{
  \begin{tabular}{c@{\hspace{3em}}c}
  Jonah Femiani & Ryan \textcolor{red}{[surname]} \\
  \textcolor{red}{[Department]} & \textcolor{red}{[Department]} \\
  \textcolor{red}{[University]} & \textcolor{red}{[University]} \\
  \textcolor{red}{[City, Country]} & \textcolor{red}{[City, Country]} \\
  \textcolor{red}{[school email]} & \textcolor{red}{[school email]}
  \end{tabular}}
\date{}

\begin{document}
\maketitle

\noindent\textbf{\textit{Abstract}---Board games give a clear, measurable setting for
testing how well a search-based agent can plan ahead. We built an agent for Connect Four
that uses depth-limited minimax search with a heuristic evaluation function and optional
alpha-beta pruning. We compared it against two baseline agents, one that plays randomly and
one that greedily takes immediate wins and blocks immediate losses, at search depths from
one to six moves. Across \X{} games per matchup, the minimax agent won \X{} percent of games
against the random agent and \X{} percent against the greedy agent at depth \X{}. Alpha-beta
pruning reduced the number of positions searched by \X{} percent at depth six without
changing the moves chosen.} \todo{state the main takeaway; rewrite once results are in.
No symbols, footnotes, or math in the title or abstract.}

\medskip
\noindent\textbf{\textit{Keywords}---game playing, adversarial search, minimax, alpha-beta
pruning, Connect Four}

\section{Introduction}

Connect Four is a two-player game played on a board seven columns wide and six rows tall.
Players take turns dropping pieces into a column, and the first to line up four pieces
horizontally, vertically, or diagonally wins. The game is small enough to program in an
afternoon but has far too many positions to solve by brute force during play, which makes it
a standard setting for studying adversarial search~[1]. With perfect play, the first player
can always force a win~[2], but finding those moves in real time takes search.

Our goal was to measure how much a classic AI method, minimax search, improves play over
simple strategies, and what that improvement costs in computation. Specifically, we asked:
\begin{enumerate}[leftmargin=*,nosep]
  \item How does the win rate of a minimax agent change as its search depth increases,
        against a random opponent and against a greedy opponent?
  \item How much does alpha-beta pruning~[3] reduce the number of positions searched and
        the time per move?
\end{enumerate}
We implemented the game, two baseline agents, and a minimax agent in Python, then ran \X{}
games per matchup with fixed random seeds. \todo{one or two sentences summarizing the main
findings.}

\section{Methodology}

\subsection{Data and Preprocessing}

This project does not use an external dataset. Instead, the data are the games played
between agents. Each matchup between two agents consists of \X{} games. The agents alternate
who moves first, because the first player has an advantage in Connect Four~[2]. Every game
uses a fixed random seed, so any game can be replayed exactly by rerunning our scripts.

For each move we record the agent, the search depth, the number of positions searched, and
the time taken. For each game we record the winner, or a draw. There is no training step, so
no train and test split is needed and there is no risk of data leakage. The minimax agent's
evaluation function was designed by hand and was not tuned on the games we report.
\todo{confirm the number of games and the seeds.}

\subsection{AI Model and Implementation}

\textit{Game representation.} The board is stored as a six-by-seven grid in which each cell
is empty or holds one player's piece. After every move, the program checks only the lines
passing through the newly placed piece, since a new win must include that piece.

\textit{Baseline agents.} The \emph{random agent} chooses uniformly among the legal
columns. The \emph{greedy agent} plays a winning move if one exists, otherwise blocks the
opponent's immediate win if one exists, and otherwise plays randomly.

\textit{Minimax agent.} Minimax~[1] explores the tree of possible future moves to a fixed
depth. It assumes the agent picks the move with the highest score and the opponent picks the
move with the lowest. Won and lost positions get very large positive or negative scores.
Positions at the depth limit are scored by a heuristic evaluation function.
\todo{Jonah: describe the final heuristic. The plan is to count the agent's open two- and
three-in-a-rows minus the opponent's, plus a bonus for center-column pieces.}

\textit{Alpha-beta pruning.} Alpha-beta pruning~[3] skips branches of the search tree that
cannot affect the final decision. It always returns the same move as plain minimax but
usually searches far fewer positions. Our implementation can switch pruning on or off, so
the two versions can be compared directly. \todo{mention move ordering if used.}

\textit{Tools.} The project is written in Python 3~[4]. The game and agents use only the
standard library, and figures are made with Matplotlib~[5]. Code is available at
\url{https://github.com/jfemiani10/connect4-bot}.

\textit{Use of generative AI.} We used Claude Code (Anthropic) to help write the game board,
terminal display, random agent, unit tests, and an early outline of this paper. We reviewed,
tested, and edited all generated code and text. \todo{update to match what was actually used,
following the course rules.}

\subsection{Experimental Setup}

We ran three experiments:
\begin{enumerate}[leftmargin=*,nosep]
  \item \emph{Win rate versus depth.} The minimax agent with alpha-beta pruning, at depths 1
        through 6, played \X{} games against each baseline.
  \item \emph{Baseline check.} The greedy agent played \X{} games against the random agent,
        to confirm that it is the stronger baseline.
  \item \emph{Cost of search.} At each depth, minimax with and without pruning was run on
        the same positions, recording positions searched and time per move.
\end{enumerate}

\textit{Metrics.} \emph{Win rate} is the fraction of games won, with draws reported
separately, and measures playing strength. \emph{Positions searched per move} counts the
board positions evaluated before choosing a move; it measures computational effort and does
not depend on hardware. \emph{Time per move} is the average number of seconds per move; it
measures practical cost but depends on hardware.

\textit{Hardware.} \todo{CPU model, RAM, operating system, and Python version used for the
timing results.}

\section{Experimental Results}

\begin{table}[t]
\caption{Win, draw, and loss rates (\%) of the minimax agent by search depth}
\label{tab:winrates}
\centering\footnotesize
\begin{tabular}{ccc}
\toprule
Depth & vs.\ Random (W / D / L) & vs.\ Greedy (W / D / L) \\
\midrule
1 & \X{} / \X{} / \X{} & \X{} / \X{} / \X{} \\
2 & \X{} / \X{} / \X{} & \X{} / \X{} / \X{} \\
3 & \X{} / \X{} / \X{} & \X{} / \X{} / \X{} \\
4 & \X{} / \X{} / \X{} & \X{} / \X{} / \X{} \\
5 & \X{} / \X{} / \X{} & \X{} / \X{} / \X{} \\
6 & \X{} / \X{} / \X{} & \X{} / \X{} / \X{} \\
\bottomrule
\end{tabular}
\end{table}

\begin{figure}[t]
\centering\figbox{figures/win\_rate\_vs\_depth.png}
\caption{Win rate (\%) of the minimax agent versus search depth, against the random and
greedy agents.}
\label{fig:winrate}
\end{figure}

\begin{figure}[t]
\centering\figbox{figures/positions\_vs\_depth.png}
\caption{Average positions searched per move versus search depth, with and without
alpha-beta pruning (log scale).}
\label{fig:positions}
\end{figure}

\begin{figure}[t]
\centering\figbox{figures/time\_vs\_depth.png}
\caption{Average time per move (s) versus search depth, with and without alpha-beta
pruning.}
\label{fig:time}
\end{figure}

Table~\ref{tab:winrates} and Fig.~\ref{fig:winrate} show the minimax agent's results against
both baselines. Fig.~\ref{fig:positions} and Fig.~\ref{fig:time} show the cost of search
with and without pruning. The greedy agent won \X{}\% of its games against the random agent.

\todo{Ryan: two or three paragraphs discussing the results. At what depth does minimax start
reliably beating each baseline? How fast do positions searched grow with depth, and how much
does pruning cut that growth? Does time per move follow the same pattern? Explain anything
unexpected, such as draws or losses at low depth.}

\section{Conclusion}

We built a Connect Four agent based on minimax search with alpha-beta pruning and compared
it with random and greedy baselines. \todo{main result about win rate versus depth; main
result about the savings from pruning. No claims not supported earlier.}

\textit{Limitations.} We only compared against weak, hand-written opponents, not strong or
human players. The evaluation function was designed by hand and was not tuned. Timing results
depend on our hardware and on Python's speed.

\textit{Next steps.} Possible extensions include testing against a perfect-play solver,
adding move ordering or a transposition table to search deeper, and comparing against a
different method such as Monte Carlo Tree Search.

\section*{References}
\todo{keep only sources we actually read; check every detail.}
\begin{enumerate}[label={[\arabic*]},leftmargin=*,itemsep=2pt,font=\footnotesize]
\footnotesize
\item S. Russell and P. Norvig, \textit{Artificial Intelligence: A Modern Approach}, 4th ed.
  Hoboken, NJ: Pearson, 2020.
\item L. V. Allis, ``A knowledge-based approach of Connect-Four: The game is solved: White
  wins,'' M.S. thesis, Vrije Universiteit, Amsterdam, The Netherlands, 1988.
\item D. E. Knuth and R. W. Moore, ``An analysis of alpha-beta pruning,'' \textit{Artificial
  Intelligence}, vol. 6, no. 4, pp. 293--326, 1975, doi: 10.1016/0004-3702(75)90019-3.
\item Python Software Foundation, ``Python 3 documentation.'' [Online]. Available:
  \url{https://docs.python.org/3/}
\item J. D. Hunter, ``Matplotlib: A 2D graphics environment,'' \textit{Computing in Science
  \& Engineering}, vol. 9, no. 3, pp. 90--95, 2007, doi: 10.1109/MCSE.2007.55.
\end{enumerate}

% ---------------------------------------------------------------------------
\onecolumn
\section*{Checklist: What's Left Before Submission}
\newcommand{\task}{\item[$\square$]}

\subsection*{Code and experiments (see TODO.md for who does what)}
\begin{itemize}[nosep]
  \task Greedy bot (Ryan).
  \task Minimax bot with heuristic, alpha-beta on/off flag, and a positions-searched counter (Jonah).
  \task \texttt{experiments/run\_matches.py} writing to \texttt{results/*.csv} (Ryan).
  \task \texttt{experiments/plots.py} producing Figs.~1--3 in \texttt{figures/} (Ryan).
  \task Run all experiments and record the hardware used.
\end{itemize}

\subsection*{Paper content}
\begin{itemize}[nosep]
  \task Fill every \X{} from the results files.
  \task Describe the final heuristic and any move ordering (Section II-B, Jonah).
  \task Write the results discussion (Section III, Ryan).
  \task Rewrite the Abstract and Introduction summary with real findings (Jonah).
  \task Rewrite the Conclusion with real findings (Ryan).
  \task Finalize the generative AI statement, following the course rules.
  \task Fill in the author details (department, university, city, email).
\end{itemize}

\subsection*{Formatting and checks}
\begin{itemize}[nosep]
  \task Copy everything into \texttt{project-template-letter.docx} using the template's
        styles. Remove the unused author blocks and the template's sample references.
  \task No symbols, special characters, footnotes, or math in the title or abstract.
  \task Every table and figure is numbered, captioned, and discussed in the text, with
        units labeled.
  \task Every reference is cited in the text, and every citation is in the reference list.
        Check the reference details.
  \task Abbreviations are defined on first use.
  \task Each teammate proofreads the other's sections.
  \task A fresh clone plus \texttt{pip install -r requirements.txt} reproduces every figure.
\end{itemize}

\end{document}
